共詞分析與文獻式知識發現係從共現（co-occurrence）解讀結構，而 AI 功能與決策功能之間稀疏的共現，很容易被解讀為人工智慧（AI）與管理決策之間存在結構洞（structural hole）。本研究所屬的專案曾作出此一推論，並據以建構多層規劃（MLP）模型。本研究說明為何共現無法承載此一構念，並以流動資料檢驗此一前提。研究一：16 個偏低的功能配對在不同隨機種子與時期之間維持其方向（在排除同詞兩端共用紀錄的規則式關鍵詞矩陣中為 13 個），但在實際可能的誤分類水準下即失去顯著性；此方向僅屬描述性。研究二對 8,512,829 件美國發明專利進行八項預先登錄的分析（其中僅一項為驗證性分析），另加一項探索性描述。AI 專利與決策專利相互引用的程度遠高於依時間分層的期望值（OE 3.05，95\% CI 2.94--3.17）；兩項相對的結構洞判準在全部六項敏感度分析中皆未成立，在移除所有帶有 G06Q10/04 分類碼的專利後亦未成立（OE 3.16）。在 16,703 組專利技術領域配對中，共現缺口同時為流動缺口的比例為 66.8\%（隨機基準為 25\%），另有 7.0\% 的流動高於期望值；以分類距離解釋兩者不一致的預先登錄假說則未獲支持。專門型 AI 公司與專門型決策公司之間的分隔，超出保留專利組合之虛無模型所能解釋的範圍（$p = .001$）：該虛無模型重現了與均等（parity）之差距的 85--91\%，但專門型公司之間的連結仍比其期望值低 22--33\%。由於沒有任何有效的虛無模型能同時固定公司度數，此一分隔尚未在固定度數下受檢驗，公司層次的中介亦仍未解決。一項預先登錄、以模型編碼的系統性文獻回顧涵蓋 204 篇已索引（indexed）研究，在其摘要中發現：135 篇提出結構洞或中介者主張，其中 73 篇建立在專利網路上；而 66 項缺口主張中有 46 項未報告任何基準。未能發現知識流動中的結構洞，並不足以確立組織層面瓶頸的存在；一項規劃中的企業研究將檢驗各項競爭解釋。
